Introduction

(TR)

Günümüzde haber üretimi ve tüketimi hızlı, çok dilli ve birçok farklı kaynağa bağlı hâle gelmiştir. Kullanıcılar yalnızca kendi ülkelerindeki haberleri değil, başka dillerde yayımlanan uluslararası gelişmeleri de takip etmek istemektedir. Ancak kaynak sayısının fazla olması, aynı konuların tekrar etmesi, makalelerin uzunluğu ve yabancı içeriklere erişimin zorluğu, açık bilgiye hızlı biçimde ulaşmayı zorlaştırmaktadır. Otomatik haber özetleme, uzun bir makaleden daha kısa, okunabilir ve bilgilendirici bir metin üreterek bu ihtiyaca cevap vermektedir. Bu proje, çok dilli haber makalelerini otomatik olarak özetleyen bütüncül bir sistem geliştirmeyi, proje kapsamında eğitilen modelleri benzer bir çalışmada eğitilmiş hazır bir modelle karşılaştırmayı ve bu modelleri gerçek zamanlı haber akışı sunan bir web uygulamasına entegre etmeyi amaçlamaktadır.

Bu bitirme projesinin temel amacı, çok dilli özetleme modellerini değerlendirmek, performanslarını karşılaştırmak ve bunları işlevsel bir web uygulamasında kullanılabilir hâle getirmektir. Proje yalnızca bir modelin eğitilmesiyle sınırlı değildir; RSS kaynaklarından makalelerin toplanmasını, metnin çıkarılmasını ve temizlenmesini, özetlerin üretilmesini, sonuçların bir veri tabanına kaydedilmesini ve filtrelenebilir bir arayüzde gösterilmesini de kapsamaktadır. Bu çerçevede, proje kapsamında mBART tabanlı modeller eğitilmiştir. mT5 modeli ise bu projede eğitilmemiştir; haber özetleme üzerine yapılmış benzer bir çalışmada önceden eğitilmiş hazır bir model olarak alınmış ve yalnızca karşılaştırma için referans model olarak kullanılmıştır.

Projenin problemi şu şekilde ifade edilebilir: Çok dilli makaleler için otomatik özetleme modelleri kalite, kaynağa bağlılık, hız ve pratik kullanılabilirlik açısından nasıl değerlendirilebilir ve ardından gerçek bir web uygulamasına nasıl entegre edilebilir? Bu soru iki boyut içermektedir. Birinci boyut modellerin performansıyla ilgilidir; çünkü farklı mimariler ve eğitim stratejileri farklı sonuçlar üretebilir. İkinci boyut ise sistemin gerçek kullanımına ilişkindir; çünkü akademik bir veri kümesinde başarılı olan bir model, gerçek haber makalelerinde farklı biçimler, gürültülü içerikler ve sürekli güncellenen kaynaklar nedeniyle çeşitli zorluklarla karşılaşabilir.

Projenin özgün değeri; model eğitimi, deneysel değerlendirme ve pratik uygulamayı bir araya getirmesinden kaynaklanmaktadır. Birçok çalışmada özetleme modelleri yalnızca sabit veri kümeleri üzerinde değerlendirilmektedir. Bu projede ise modeller, RSS akışlarından gelen güncel haberlerle çalışan bir sisteme de entegre edilmiştir. Uygulama haberleri toplamakta, mümkün olduğunda ana metni çıkarmakta, bu metni temizlemekte ve seçilen modelle özetler üretmektedir. Böylece proje yalnızca teorik bir karşılaştırma sunmakla kalmamakta, gerçek kullanım senaryosuna yakın bir prototip de ortaya koymaktadır.

Bir diğer özgün yön, yeni proxy skorlarla çok boyutlu bir değerlendirme tasarlanmasıdır. ROUGE, BLEU veya METEOR gibi klasik metrikler yararlıdır; ancak bir haber özetinin kalitesini tamamen değerlendirmek için tek başına yeterli değildir. İyi bir özet aynı zamanda kaynak metne sadık kalmalı, tekrarları önlemeli, açık olmalı, önemli bilgileri kapsamalı ve makul bir sürede üretilmelidir. Bu nedenle projede kaynak kapsamı, kaynak geri çağırımı, tekrar oranı, sıkıştırma oranı ve üretim süresi gibi ölçümler de kullanılmaktadır. Bu sinyallerden hareketle tutarlılık, doğruluk, açıklık, ilgililik ve verimlilik skorları oluşturulmuştur. Bu yaklaşım yalnızca “hangi model klasik metriklerde en iyi sonucu alıyor?” sorusuna değil, aynı zamanda “hangi model hangi kullanım amacı için daha uygundur?” sorusuna da cevap vermeyi sağlamaktadır.

Kullanılan yöntemler üç ana aşamada ele alınabilir. Birinci aşama modellerin hazırlanması ve eğitilmesidir. BART, mBART ve mT5 tabanlı sequence-to-sequence transformer modelleri incelenmiştir. BART modelleri İngilizce özetleme için değerlendirilmiş, mBART modelleri ise çok dilli özetleme için kullanılmıştır. mBART için veri kümesi, hedef özet uzunluğu, öğrenme oranı, optimizasyon, doğrulama ve en iyi checkpoint seçimi gibi çeşitli eğitim tercihleri test edilmiştir. mT5 modeli bu eğitimin bir sonucu değildir; benzer bir çalışmada eğitilmiş hazır bir model olarak alınmış ve karşılaştırma noktası olarak kullanılmıştır.

İkinci aşama modellerin değerlendirilmesi ve karşılaştırılmasıdır. Projede eğitilen mBART modelleri ile hazır mT5 modeli aynı koşullar altında test edilmiştir. Bu nedenle mT5 sonuçları, projenin eğitim çıktısı olarak değil, mBART modellerini hâlihazırda güçlü olan bir modelle konumlandırmaya yarayan bir referans olarak değerlendirilmelidir. Bu değerlendirme sırasında klasik metrikler, bileşik proxy skorlarla tamamlanmıştır. Bu skorlar, klasik metriklerin tek başına yakalayamadığı boyutları ölçmek amacıyla tasarlanmıştır.

Bu çerçevede tutarlılık, doğruluk, açıklık, ilgililik ve verimlilik skorları ayrı ayrı hesaplanmış, ardından “capability overall” adı verilen genel bir değer altında birleştirilmiştir. “Quality factuality” skoru özetin kaynak metne sadakatini ölçmek için, “quality completeness” skoru ise referans içeriğin kapsanmasını ve özet uzunluğunun uygunluğunu değerlendirmek için kullanılmıştır. Bu skorlar insan değerlendirmesinin yerini almaz; ancak modellerin otomatik, tekrarlanabilir ve daha ayrıntılı biçimde karşılaştırılmasına olanak sağlar.

Üçüncü aşama modellerin web uygulamasına entegre edilmesidir. Flask tabanlı bir uygulama geliştirilmiştir. Bu uygulama farklı dil ve kategoriler için RSS akışlarını toplamakta, makalelerin ana metnini çıkarmakta, içerikleri temizlemekte ve metinleri özetleme modellerine göndermektedir. Üretilen özetler; makale, kaynak, dil, model ve tarih bilgileriyle birlikte SQLite veri tabanına kaydedilmektedir. Kullanıcı arayüzü daha sonra önceden hazırlanmış özetleri okumaktadır; böylece her kullanıcı isteğinde modelin yeniden çalıştırılmasına gerek kalmamaktadır. Uygulama; dil, konu, bölge, ülke, kaynak, anahtar kelime ve model bazında filtreleme seçeneklerinin yanı sıra çeviri ve görsel gösterimi gibi özellikler de sunmaktadır.

Sistemin tasarımı aynı zamanda arka planda çalışan bir veri alma mimarisine dayanmaktadır. Makaleleri tam olarak kullanıcı isteği anında toplamak ve özetlemek yerine, sistem haberleri düzenli aralıklarla bir zamanlayıcı yardımıyla işlemektedir. Bu yapı, RSS akışlarından aday makaleleri toplamakta, kaynaklar ve diller arasında işlem yükünü dengelemekte, metinleri çıkarmakta, özetler üretmekte ve bunları veri tabanına kaydetmektedir. Bu yaklaşım kullanıcı bekleme süresini azaltmakta ve uygulamayı daha kararlı hâle getirmektedir.

Projenin bilimsel katkısı, çok dilli özetleme modellerini tek bir performans metriğinin ötesine geçen daha geniş bir çerçevede değerlendirmesidir. Yeni proxy skorlar; kaliteyi tutarlılık, doğruluk, açıklık, ilgililik, verimlilik, kaynağa sadakat ve içerik bütünlüğü gibi farklı açılardan analiz etmeyi mümkün kılmaktadır. Çalışma ayrıca farklı eğitim stratejilerinin mBART modellerinin davranışını nasıl etkilediğini gözlemlemeye olanak tanımaktadır. Benzer bir çalışmada eğitilmiş mT5 modelinin eklenmesi güçlü bir karşılaştırma noktası sunmakta ve analizi daha anlamlı hâle getirmektedir.

Teknolojik açıdan proje, yapay zekâ modellerinin gerçek bir yazılım sistemine nasıl entegre edilebileceğini göstermektedir. Proje; yerel model kullanımını, RSS toplama sürecini, makale çıkarımını, arka plan işlemeyi, veri tabanı depolamayı ve web arayüzünü birleştirmektedir. Bu mimari, başka özetleme veya otomatik içerik işleme uygulamaları için de temel oluşturabilir. Aynı makaleye birden fazla modelin uygulanabilmesi, model çıktılarının pratik bir bağlamda karşılaştırılmasını da mümkün kılmaktadır.

Etik açıdan proje, otomatik özetlemenin sınırlılıklarını dikkate almaktadır. Bir model bazı ayrıntıları çıkarabilir, içeriği aşırı basitleştirebilir veya eksik bilgi üretebilir. Bu nedenle uygulamada orijinal makaleye giden bağlantı her zaman korunmaktadır. Özetler, makalenin yerini alan nihai bir gerçeklik olarak değil, hızlı okumaya yardımcı olan bir araç olarak sunulmaktadır. Ayrıca kaynak metne sadakatle ilgili metriklerin kullanılması, orijinal içerikten fazla uzaklaşan özetlerin riskini azaltmaya katkı sağlamaktadır.

Sosyoekonomik açıdan proje, bilgiye erişimi kolaylaştırabilir. Uzun makaleleri okumaya zamanı olmayan kullanıcılar, özetler sayesinde gündemi daha hızlı takip edebilir. Çok dilli destek, yabancı kaynakları daha erişilebilir hâle getirir. Konu, ülke, bölge ve kaynak bazında filtreleme; öğrenciler, araştırmacılar, uluslararası gündemi takip eden kullanıcılar ve medya izleme alanında çalışan kurumlar için yararlı olabilir. Böyle bir sistem, bilgiye erişim maliyetini azaltabilir ve dijital haber tüketiminin verimliliğini artırabilir.

Projenin temel çıktıları üç kategoride toplanmaktadır. Birinci çıktı, proje kapsamında eğitilen çok dilli mBART modellerinin oluşturulması ve kullanılmasıdır. Hazır mT5 modeli, projenin eğitim çıktısı değildir; yalnızca karşılaştırma amacıyla eklenmiştir. İkinci çıktı, modellerin klasik metrikler ve proxy skorlarla karşılaştırmalı olarak değerlendirilmesidir. Üçüncü çıktı ise haberleri toplayan, özetler üreten, bunları veri tabanında saklayan ve kullanıcıya sunan işlevsel bir web uygulamasıdır.

Sonuçlar, hazır mT5 modelinin güçlü bir karşılaştırma noktası oluşturduğunu; proje kapsamında eğitilen mBART modellerinin ise farklı senaryolar için faydalı profiller sunduğunu göstermektedir. mT5’in bazı ölçütlerde daha iyi skorlar alması, proje modellerinin başarısızlığı olarak yorumlanmamalıdır; bu durum, zaten eğitilmiş ve güçlü bir modelle yapılan karşılaştırmanın sonucudur. mBART modelleri arasında seçim, önceliklere bağlıdır: kaynak metne sadakat, hız, referansla benzerlik veya genel denge. Bu nedenle proje, model seçiminin hedeflenen kullanım amacına göre yapılması gerektiğini göstermektedir.

Çalışma planı aşamalı olarak gerçekleştirilmiştir. İlk olarak problem tanımlanmış ve projenin kapsamı çok dilli özetleme, model karşılaştırması ve web entegrasyonu ile sınırlandırılmıştır. Daha sonra veri kümeleri, RSS kaynakları ve BART, mBART ve mT5 tabanlı aday modeller incelenmiştir. Ardından mBART modelleri farklı parametrelerle eğitilmiş ve yerel olarak kullanılabilir hâle getirilmiştir; benzer bir çalışmada önceden eğitilmiş olan mT5 modeli ise karşılaştırma modeli olarak eklenmiştir. Bundan sonra klasik metrikler, kaynak sadakati, okunabilirlik, tekrar, sıkıştırma, üretim süresi ve yeni proxy skorları içeren değerlendirme süreci tasarlanmıştır. Modeller aynı koşullar altında test edilmiş ve sonuçları karşılaştırılmıştır. Bir sonraki aşamada RSS havuzu, metin çıkarımı, temizleme, özet üretimi, SQLite, ingestion zamanlayıcısı ve API endpointleriyle birlikte web uygulaması geliştirilmiştir. Son olarak kullanıcı arayüzü hazırlanmış, sistem uçtan uca test edilmiş ve deneysel sonuçlar ile bilimsel, teknolojik, etik ve sosyoekonomik katkılar analiz edilmiştir. 


---

(FR)

Aujourd’hui, la production et la consommation d’actualités sont devenues rapides, multilingues et liées à de nombreuses sources. Les utilisateurs veulent suivre non seulement les nouvelles de leur propre pays, mais aussi les événements internationaux publiés dans d’autres langues. Cependant, le grand nombre de sources, la répétition des mêmes sujets, la longueur des articles et la difficulté d’accéder aux contenus étrangers rendent l’accès rapide à une information claire plus difficile. Le résumé automatique de nouvelles répond à ce besoin en produisant un texte plus court, lisible et informatif à partir d’un article long. Ce projet vise à développer un système complet qui résume automatiquement des articles multilingues, compare les modèles entraînés dans le projet avec un modèle prêt à l’emploi entraîné dans une étude similaire, et intègre ces modèles dans une application web de flux d’actualités en temps réel.

L’objectif principal de ce projet de fin d’études est d’évaluer des modèles de résumé multilingue, de comparer leurs performances et de les rendre utilisables dans une application web fonctionnelle. Le projet ne se limite pas à l’entraînement d’un modèle : il comprend aussi la collecte d’articles depuis des sources RSS, l’extraction et le nettoyage du texte, la génération de résumés, l’enregistrement des résultats dans une base de données et leur affichage dans une interface filtrable. Dans ce cadre, les modèles basés sur mBART ont été entraînés dans le projet. Le modèle mT5, lui, n’a pas été entraîné dans ce projet ; il a été pris comme modèle déjà entraîné dans une étude similaire sur le résumé d’actualités et utilisé uniquement comme modèle de référence pour la comparaison.

La problématique du projet peut être formulée ainsi : comment évaluer des modèles automatiques de résumé pour des articles multilingues selon la qualité, la fidélité à la source, la vitesse et l’utilisabilité pratique, puis comment les intégrer dans une application web réelle ? Cette question contient deux dimensions. La première concerne la performance des modèles, car différentes architectures et stratégies d’entraînement peuvent produire des résultats différents. La deuxième concerne l’utilisation réelle du système, car un modèle performant sur un jeu de données académique peut rencontrer des difficultés avec des articles réels, souvent bruités, de formats variés et constamment mis à jour.

La valeur originale du projet vient de l’union entre entraînement de modèles, évaluation expérimentale et application pratique. Dans beaucoup d’études, les modèles de résumé sont évalués seulement sur des jeux de données fixes. Ici, les modèles sont aussi intégrés dans un système qui travaille avec des actualités actuelles provenant de flux RSS. L’application collecte les nouvelles, extrait si possible le texte principal, nettoie ce texte et produit des résumés avec le modèle choisi. Ainsi, le projet ne propose pas seulement une comparaison théorique, mais aussi un prototype proche d’un vrai scénario d’utilisation.

Un autre aspect original est la conception d’une évaluation multidimensionnelle avec de nouveaux scores proxy. Les métriques classiques comme ROUGE, BLEU ou METEOR sont utiles, mais elles ne suffisent pas pour juger complètement la qualité d’un résumé d’actualité. Un bon résumé doit aussi rester fidèle au texte source, éviter les répétitions, être clair, couvrir l’information importante et être produit dans un temps raisonnable. Pour cette raison, le projet utilise aussi des mesures comme la couverture de la source, le rappel de la source, le taux de répétition, le taux de compression et le temps de génération. À partir de ces signaux, des scores de cohérence, d’exactitude, de clarté, de pertinence et d’efficacité ont été créés. Cette approche permet de répondre non seulement à la question « quel modèle obtient le meilleur score classique ? », mais aussi à la question « quel modèle convient le mieux à tel usage ? ».

Les méthodes utilisées peuvent être regroupées en trois étapes principales. La première étape est la préparation et l’entraînement des modèles. Des modèles transformer de type séquence-à-séquence basés sur BART, mBART et mT5 ont été étudiés. Les modèles BART ont été considérés pour le résumé en anglais, tandis que les modèles mBART ont été utilisés pour le résumé multilingue. Plusieurs choix d’entraînement ont été testés pour mBART, comme le jeu de données, la longueur cible du résumé, le taux d’apprentissage, l’optimisation, la validation et le choix du meilleur checkpoint. Le modèle mT5 n’est pas un résultat de cet entraînement ; il a été pris comme modèle prêt à l’emploi entraîné dans une étude similaire et utilisé comme point de comparaison.

La deuxième étape est l’évaluation et la comparaison des modèles. Les modèles mBART entraînés dans le projet et le modèle mT5 prêt à l’emploi ont été testés dans les mêmes conditions. Les résultats de mT5 ne représentent donc pas un résultat d’entraînement du projet, mais une référence permettant de situer les modèles mBART face à un modèle déjà fort. Pendant cette évaluation, les métriques classiques ont été complétées par des scores proxy composites. Ces scores ont été conçus pour mesurer des dimensions que les métriques classiques ne capturent pas seules.

Dans ce cadre, les scores de cohérence, d’exactitude, de clarté, de pertinence et d’efficacité ont été calculés séparément, puis combinés dans une valeur générale appelée capability overall. Le score quality factuality a été utilisé pour mesurer la fidélité du résumé au texte source, tandis que quality completeness a servi à évaluer la couverture du contenu de référence et l’adaptation de la longueur du résumé. Ces scores ne remplacent pas une évaluation humaine, mais ils permettent une comparaison automatique, répétable et plus détaillée des modèles.

La troisième étape est l’intégration des modèles dans l’application web. Une application basée sur Flask a été développée. Elle collecte des flux RSS pour différentes langues et catégories, extrait le texte principal des articles, nettoie les contenus et envoie les textes aux modèles de résumé. Les résumés produits sont enregistrés dans une base de données SQLite avec les informations sur l’article, la source, la langue, le modèle et la date. L’interface utilisateur lit ensuite les résumés déjà préparés, ce qui évite de relancer le modèle à chaque requête. L’application propose des filtres par langue, sujet, région, pays, source, mot-clé et modèle, ainsi que des options comme la traduction et l’affichage d’images.

La conception du système repose aussi sur une architecture d’ingestion en arrière-plan. Au lieu de collecter et résumer les articles au moment exact de la demande utilisateur, le système traite régulièrement les nouvelles avec un planificateur. Celui-ci collecte les articles candidats depuis les flux RSS, équilibre le traitement entre les sources et les langues, extrait les textes, produit les résumés et les enregistre dans la base de données. Cette approche réduit le temps d’attente pour l’utilisateur et rend l’application plus stable.

La contribution scientifique du projet est l’évaluation des modèles de résumé multilingue avec un cadre plus large qu’une seule métrique de performance. Les nouveaux scores proxy permettent d’analyser la qualité selon plusieurs aspects : cohérence, exactitude, clarté, pertinence, efficacité, fidélité à la source et complétude du contenu. L’étude permet aussi d’observer comment différentes stratégies d’entraînement influencent le comportement de modèles mBART. L’ajout du modèle mT5 entraîné dans une étude similaire donne un point de comparaison fort et rend l’analyse plus significative.

Du point de vue technologique, le projet montre comment intégrer des modèles d’intelligence artificielle dans un vrai système logiciel. Il combine l’utilisation locale de modèles, la collecte RSS, l’extraction d’articles, le traitement en arrière-plan, le stockage en base de données et une interface web. Cette architecture peut servir de base pour d’autres applications de résumé ou de traitement automatique de contenu. Le fait de pouvoir appliquer plusieurs modèles au même article permet aussi de comparer leurs sorties dans un contexte pratique.

Du point de vue éthique, le projet prend en compte les limites du résumé automatique. Un modèle peut supprimer certains détails, simplifier excessivement le contenu ou produire une information incomplète. Pour cette raison, le lien vers l’article original est toujours conservé dans l’application. Les résumés ne sont pas présentés comme une vérité finale remplaçant l’article, mais comme une aide à la lecture rapide. De plus, l’utilisation de métriques liées à la fidélité au texte source contribue à réduire le risque de résumés trop éloignés du contenu original.

Du point de vue socioéconomique, le projet peut faciliter l’accès à l’information. Les utilisateurs qui n’ont pas le temps de lire de longs articles peuvent suivre plus rapidement l’actualité grâce aux résumés. Le support multilingue rend les sources étrangères plus accessibles. Le filtrage par sujet, pays, région et source peut être utile pour les étudiants, les chercheurs, les utilisateurs suivant l’actualité internationale et les institutions travaillant sur le suivi des médias. Un tel système peut donc réduire le coût d’accès à l’information et améliorer l’efficacité de la consommation d’actualités numériques.

Les principaux résultats du projet sont regroupés en trois catégories. Le premier résultat est la création et l’utilisation de modèles mBART multilingues entraînés dans le cadre du projet. Le modèle mT5 prêt à l’emploi n’est pas un résultat d’entraînement du projet ; il a été ajouté uniquement pour la comparaison. Le deuxième résultat est l’évaluation comparative des modèles avec des métriques classiques et des scores proxy. Le troisième résultat est une application web fonctionnelle qui collecte les actualités, produit les résumés, les stocke dans une base de données et les présente à l’utilisateur.

Les résultats montrent que le modèle mT5 prêt à l’emploi constitue un point de comparaison fort, tandis que les modèles mBART entraînés dans le projet produisent des profils utiles pour différents scénarios. Le fait que mT5 obtienne de meilleurs scores sur certaines mesures ne doit pas être interprété comme un échec des modèles du projet, mais comme une comparaison avec un modèle déjà entraîné et solide. Entre les modèles mBART, le choix dépend des priorités : fidélité à la source, vitesse, similarité avec la référence ou équilibre général. Le projet montre donc que le choix du modèle doit dépendre de l’usage visé.

Le plan de travail a été réalisé par étapes. D’abord, la problématique a été définie et le périmètre du projet a été limité au résumé multilingue, à la comparaison de modèles et à l’intégration web. Ensuite, les jeux de données, les sources RSS et les modèles candidats basés sur BART, mBART et mT5 ont été étudiés. Puis, les modèles mBART ont été entraînés avec différents paramètres et rendus utilisables localement ; le modèle mT5, déjà entraîné dans une étude similaire, a été ajouté comme modèle de comparaison. Après cela, le processus d’évaluation a été conçu avec les métriques classiques, les métriques de fidélité à la source, de lisibilité, de répétition, de compression, de temps de génération et les nouveaux scores proxy. Les modèles ont ensuite été testés dans les mêmes conditions et leurs résultats ont été comparés. Dans l’étape suivante, l’application web a été développée avec le pool RSS, l’extraction des textes, le nettoyage, la génération de résumés, SQLite, le planificateur d’ingestion et les endpoints API. Enfin, l’interface utilisateur a été préparée, le système a été testé de bout en bout, puis les résultats expérimentaux et les contributions scientifiques, technologiques, éthiques et socioéconomiques ont été analysés.